\section{The equal and proportional cases}\label{sec:ft_equal_proportional_cases}

% The CPSV16 proportional theorem is proved here for BNT canonical forms,
% using the sector-matching witnesses of Chapter~\ref{ch:bnt}. The
% per-block phases are not collapsed into a single global gauge on the total
% tensor; that additional step needs the proportional coefficient identity and
% is recorded conditionally above
% (thm:sector_bnt_proportional_global_gauge_of_coeff_identity).
\begin{theorem}[Multi-block fundamental theorem for normal tensors]%
    \label{thm:cpgsv_multiblock_ft_source}
    \lean{MPSTensor.fundamentalTheorem_proportional_canonicalForm}
    \leanok
    \uses{def:sector_bnt_cf, def:gauge_phase_equiv,
        thm:sector_bnt_proportional_sector_match_witnesses}
    Let $P$ and $Q$ be BNT canonical forms (Definition~\ref{def:sector_bnt_cf}),
    with bases of normal tensors $(P_j)_{j=1}^{g_a}$ and $(Q_k)_{k=1}^{g_b}$. If
    the total tensors $P_{\mathrm{tot}}$, $Q_{\mathrm{tot}}$ generate eventually
    nonzero proportional MPV families --- for all sufficiently large $N$ there is
    a nonzero scalar $c_N$ with
    $V^{(N)}(P_{\mathrm{tot}})=c_NV^{(N)}(Q_{\mathrm{tot}})$ --- then
    $g_a=g_b$. After relabelling the basis of normal tensors of $P$ by a
    bijection $\beta:J(Q)\simeq J(P)$, there are scalars $\zeta_k\in\C$ with
    $|\zeta_k|=1$ and invertible matrices $Y_k$ such that, for every $k$ and
    every physical index~$i$,
    \begin{align}
        Q_k^i
        &=\zeta_k\,Y_k\,P_{\beta(k)}^i\,Y_k^{-1}. \notag
    \end{align}
    Thus the paired normal tensors have the same bond dimension
    ($\dim P_{\beta(k)}=\dim Q_k$) and the same gauge-phase class.
\end{theorem}

\begin{proof}\leanok
    \uses{thm:sector_bnt_proportional_sector_match_witnesses}
    The proportional sector-matching
    Theorem~\ref{thm:sector_bnt_proportional_sector_match_witnesses}, applied to
    the BNT canonical forms $P$ and $Q$, supplies the bijection $\beta$, the
    bond-dimension equalities, the unit phases $\zeta_k$, and the block gauges
    $Y_k$ with $Q_k^i=\zeta_kY_kP_{\beta(k)}^iY_k^{-1}$; the matching needs
    no per-sector unit-modulus hypothesis.
\end{proof}

This is \cite[Theorem~2.10]{Cirac2016MPDO_arXiv}; see also
\cite[Theorem~IV.4 and Corollary~IV.5]{Cirac2021Matrix}. The statement is given on
the BNT canonical-form surface (Definition~\ref{def:sector_bnt_cf}), which is the
paper's ``$A$ and $B$ in canonical form'' with chosen bases of normal tensors;
the reduction of an \emph{arbitrary} proportional-MPV pair to a common primitive
block decomposition is
Theorem~\ref{thm:ft_after_blocking_common_length_common_sector_theorem} of
Chapter~\ref{ch:canonical}.

\begin{definition}[Left-canonical normalization]
    \label{def:left_canonical_tensor}
    \lean{MPSTensor.IsLeftCanonical}
    \leanok
    An MPS tensor $A=(A^i)_i$ is left-canonical when
    \begin{align}
        \sum_i (A^i)^\dagger A^i&=\Id. \notag
    \end{align}
\end{definition}

\begin{lemma}[Phase normalization for an irreducible gauge]
    \label{lem:left_canonical_irreducible_phase_norm}
    \lean{MPSTensor.gaugePhase_scalar_norm_eq_one_of_leftCanonical_irreducible}
    \leanok
    \uses{def:left_canonical_tensor, def:irreducible_tensor,
        thm:eigenvalue_unique_irred}
    Let $A$ and $B$ be left-canonical tensors of common positive bond
    dimension, with $B$ irreducible. If $X$ is invertible, $\zeta\ne0$, and
    $B^i=\zeta XA^iX^{-1}$ for every $i$, then $|\zeta|=1$.
\end{lemma}

\begin{proof}\leanok
    \uses{thm:eigenvalue_unique_irred}
    Let $\rho$ be a nonzero positive fixed point of $\E_A$ and set
    $\sigma=X\rho X^\dagger$. Then $\sigma$ is nonzero and positive, and
    $\E_B(\sigma)=|\zeta|^2\sigma$. The map $\E_B$ also has a nonzero positive
    fixed point. Its irreducibility therefore forces $|\zeta|^2=1$.
\end{proof}

\begin{lemma}[Unitary replacement with the same phase]
    \label{lem:left_canonical_irreducible_same_phase_unitary_gauge}
    \lean{MPSTensor.exists_unitaryConj_of_gaugePhase_data_of_leftCanonical_irreducible}
    \leanok
    \uses{def:left_canonical_tensor, def:irreducible_tensor,
        lem:left_canonical_irreducible_phase_norm, thm:psd_fp_unique_irred}
    Let $A$ and $B$ be left-canonical irreducible tensors of common positive
    bond dimension satisfying $B^i=\zeta XA^iX^{-1}$ for an invertible $X$ and
    a nonzero scalar $\zeta$. There is a unitary matrix $U$, with the same scalar
    $\zeta$, such that $|\zeta|=1$ and $B^i=\zeta UA^iU^\dagger$ for every $i$.
\end{lemma}

\begin{proof}\leanok
    \uses{lem:left_canonical_irreducible_phase_norm, thm:psd_fp_unique_irred}
    Lemma~\ref{lem:left_canonical_irreducible_phase_norm} gives
    $|\zeta|=1$. Substitution into the left-canonical identities gives
    \begin{align}
        \sum_i A^{i\dagger}(X^\dagger X)A^i&=X^\dagger X. \notag
    \end{align}
    Thus $X^\dagger X$ is a positive-semidefinite fixed point of the adjoint
    transfer map of $A$. Irreducibility gives $X^\dagger X=c\Id$ for some
    $c>0$, so $X^{-1}=c^{-1}X^\dagger$ and $U=c^{-1/2}X$ is unitary. Moreover,
    \begin{align}
        UA^iU^\dagger
        &=c^{-1}XA^iX^\dagger
        =XA^iX^{-1}. \notag
    \end{align}
\end{proof}

\begin{lemma}[Unitary gauge for left-canonical irreducible tensors]
    \label{lem:left_canonical_irreducible_unitary_gauge}
    \lean{MPSTensor.exists_unitaryConj_gaugePhase_of_leftCanonical_irreducible}
    \leanok
    \uses{def:left_canonical_tensor, def:gauge_phase_equiv,
        def:irreducible_tensor,
        lem:left_canonical_irreducible_same_phase_unitary_gauge}
    Let $A$ and $B$ be left-canonical irreducible tensors of common positive bond
    dimension that are gauge-phase equivalent. Then the gauge can be taken
    unitary: there are a unitary matrix $U$ and a scalar $\zeta$ with
    $|\zeta|=1$ such that, for every physical index~$i$,
    \begin{align}
        B^i&=\zeta\,UA^iU^\dagger. \notag
    \end{align}
\end{lemma}

\begin{proof}\leanok
    \uses{lem:left_canonical_irreducible_same_phase_unitary_gauge}
    Unpack the gauge-phase equivalence as
    $B^i=\zeta XA^iX^{-1}$ with $X$ invertible and $\zeta\ne0$.
    Lemma~\ref{lem:left_canonical_irreducible_same_phase_unitary_gauge}
    replaces $X$ by a unitary $U$ while retaining this same phase and proves
    $|\zeta|=1$.
\end{proof}

\begin{corollary}[Canonical Form II unitary refinement of proportional matching]%
    \label{cor:sector_bnt_proportional_unitary_sector_match}
    \lean{MPSTensor.ft_sector_bnt_proportional_unitary_sector_match_witnesses}
    \leanok
    \uses{def:sector_bnt_cf, thm:cpgsv_multiblock_ft_source}
    Under the hypotheses of
    Theorem~\ref{thm:cpgsv_multiblock_ft_source}, there are a bijection
    $\beta:J(Q)\simeq J(P)$, unitary matrices
    $U_k\in\mathrm{U}(\dim Q_k)$, and phases $\zeta_k\in\C$ with
    $|\zeta_k|=1$, such that $\dim P_{\beta(k)}=\dim Q_k$ and, for every $k$
    and every physical index~$i$,
    \begin{align}
        Q_k^i
        &=\zeta_k\,U_k\,P_{\beta(k)}^i\,U_k^\dagger. \notag
    \end{align}
    This is the proportional part of
    \cite[Corollary~C.5]{Cirac2016MPDO_arXiv}.
\end{corollary}

\begin{proof}\leanok
    \uses{thm:sector_bnt_proportional_sector_match_witnesses,
        lem:left_canonical_irreducible_unitary_gauge}
    The proportional sector-matching theorem gives a bijection $\beta$,
    invertible matrices $X_k$, and unit-modulus phases $\zeta_k$ with
    $Q_k^i=\zeta_kX_kP_{\beta(k)}^iX_k^{-1}$. The matched tensors are
    left-canonical and irreducible. Hence
    Lemma~\ref{lem:left_canonical_irreducible_unitary_gauge} gives unitary
    matrices $U_k$ satisfying
    $X_kP_{\beta(k)}^iX_k^{-1}=U_kP_{\beta(k)}^iU_k^\dagger$, and therefore
    $Q_k^i=\zeta_kU_kP_{\beta(k)}^iU_k^\dagger$.
\end{proof}

% In the equal case the per-block phases and copy-weight identities assemble
% into the unitary global gauge of
% cor:sector_bnt_equal_unitary_global_gauge. Under proportional input, the
% analogous total-tensor statement remains conditional on the proportional
% coefficient identity.

\begin{theorem}[Equal BNT canonical forms are globally conjugate]%
    \label{thm:cpgsv_equal_case_source}
    \lean{MPSTensor.fundamentalTheorem_equal_canonicalForm}
    \leanok
    \uses{def:sector_bnt_cf, thm:sector_bnt_equal_mps_gaugeEquiv_literal}
    Let $P$ and $Q$ be BNT canonical forms (Definition~\ref{def:sector_bnt_cf})
    whose total tensors $P_{\mathrm{tot}}$, $Q_{\mathrm{tot}}$ generate the same
    MPV family for every positive system size~$N$. Then their total bond
    dimensions coincide; writing the common value as~$D$, there is one
    invertible matrix $X\in\GL_D(\C)$ such that, for every physical index~$i$,
    \begin{align}
        Q_{\mathrm{tot}}^i
        &=X\,P_{\mathrm{tot}}^i\,X^{-1}. \notag
    \end{align}

    This is \cite[Corollary~2.11]{Cirac2016MPDO_arXiv} on the BNT
    canonical-form surface; the canonical form of
    \cite[Section~2.3]{Cirac2016MPDO_arXiv}, a direct sum of normal tensors with
    scalar weights, is a BNT canonical form. The reduction of an \emph{arbitrary}
    pair of same-MPV tensors to a common BNT canonical form is
    Chapter~\ref{ch:canonical}.
\end{theorem}

\begin{proof}\leanok
    \uses{thm:sector_bnt_equal_mps_gaugeEquiv_literal}
    This is Theorem~\ref{thm:sector_bnt_equal_mps_gaugeEquiv_literal}. The BNT
    comparison of Chapter~\ref{ch:bnt} provides the basis relabelling, the
    per-block gauge-phase equivalences, and the copy-weight multiset matching,
    which combine into the block-diagonal global gauge. The exact
    linear-independence matching needs no per-sector unit-modulus hypothesis, and
    the BNT canonical form needs no one-site injectivity.
\end{proof}

This is also \cite[Corollary~IV.5]{Cirac2021Matrix}.
